% !TeX root = ../../Book2.tex
% 纲要标题：### D.5 与群、张量和非交换几何的联系

% 纲要要点（供目录核对）：
% * 从离散群代数 \(\mathbb C[G]\) 到全群／约化群 \(C^*\)-代数及群 von Neumann 代数；以 \(G=\mathbb Z\) 比较 Laurent 多项式与圆周函数
% * 代数张量积与算子代数张量积需要区分；定义最小／最大 \(C^*\)-张量范数，引用其极值定理并说明含幺情形的最大性
% * 高阶算子 K 理论、Tomita–Takesaki 理论、非交换几何和非自伴算子代数列为专题出口；本附录不是这些专著的压缩替代；作为本节末阅读说明，不另设小节
% #### D.5.1 离散群的算子代数
% #### D.5.2 代数张量积与算子张量范数
% ---

\section{与群、张量和非交换几何的联系}
\label{b2:sec:D05}

同一个群代数可以有不同的完备化，同一个代数张量积也可以有不同的算子范数。本节只建立辨认这些对象所需的定义，并在整数群及矩阵代数上计算。这样既能看到与本卷代数构造的联系，也能说明后续分析问题究竟增加了什么。

\subsection{离散群的算子代数}
\label{b2:subsec:D05:01}

对离散群 \(G\)，群代数 \(\mathbb C[G]\) 的元素是有限和 \(a=\sum_g a_g g\)。规定
\[
a^*=\sum_g\overline{a_g}g^{-1}.
\]
任何酉表示 \(U:G\to\mathcal U(H)\) 通过线性延拓给出 \(*\)-表示 \(U(a)=\sum_g a_g U_g\)，且
\[
\bigl\|U(a)\bigr\|\le\sum_g|a_g|.
\]
所以 \(\|a\|_{\max}=\sup_U\bigl\lVert U(a)\bigr\rVert\) 有限。它不只是一项半范数：在 \(\ell^2(G)\) 的标准正交基上定义左正则表示 \(\lambda_g\delta_h=\delta_{gh}\)，则
\(\lambda(a)\delta_e=\sum_g a_g\delta_g\)，故 \(a\ne0\) 时 \(\lambda(a)\ne0\)。次乘性、对合及 \(C^*\) 恒等式由各算子表示继承，得到一个真正的 \(C^*\) 范数。

\begin{definition}\label{b2:def:group-operator-algebras}
设 \(G\) 为离散群，\(\mathbb C[G]\) 为其复群代数。对 \(a=\sum_g a_g g\in\mathbb C[G]\)，令 \(\|a\|_{\max}=\sup_U\|\sum_g a_gU_g\|\)，其中 \(U\) 遍历 \(G\) 在复 Hilbert 空间上的酉表示。将 \(\mathbb C[G]\) 按此范数完备化，所得代数称为\term[quan-qun-Cxing-daishu]{全群 \(C^*\)-代数}[full group C*-algebra]，记为 \(C^*(G)\)。在 \(\ell^2(G)\) 的标准基上令 \(\lambda_g\delta_h=\delta_{gh}\)。将左正则表示 \(\lambda:\mathbb C[G]\to B(\ell^2(G))\) 的像取算子范数闭包，所得代数称为\term[yuehua-qun-Cxing-daishu]{约化群 \(C^*\)-代数}[reduced group C*-algebra]，记为 \(C_r^*(G)\)；若取弱算子拓扑下的闭包，则称所得代数为\term[qun-von-Neumann-daishu]{群 von Neumann 代数}[group von Neumann algebra]，记为 \(L(G)\)。
\end{definition}
\symidx{\(C^*(G)\)：离散群的全群 \(C^*\)-代数}
\symidx{\(C_r^*(G)\)：离散群的约化群 \(C^*\)-代数}
\symidx{\(L(G)\)：群 von Neumann 代数}

\(L(G)=\lambda(G)''\) 由\cref{b2:thm:operator-bicommutant}成立。每个酉表示唯一延拓到 \(C^*(G)\)，因为它对最大范数收缩；但不能不加条件地断言它对约化范数也收缩。全群与约化群代数间的自然映射及其差异见 Blackadar\cite[\nopp II.10.2.3--II.10.2.7]{Blackadar2017OperatorAlgebrasCorrected}。

\begin{proposition}\label{b2:prop:integer-group-cstar}
设 \(G=\mathbb Z\)，以 \(u\) 表示群代数 \(\mathbb C[G]\) 中对应于 \(1\in G\) 的元素，并令 \(u^*=u^{-1}\)。记单位圆为 \(\mathbb T=\bigl\{\,z\in\mathbb C\,\bigm|\,|z|=1\,\bigr\}\)，\(C(\mathbb T)\) 为其上的复连续函数代数，则 \(\mathbb C[G]=\mathbb C[u,u^{-1}]\)，且
\[
C^*(\mathbb Z)\cong C_r^*(\mathbb Z)\cong C(\mathbb T).
\]
\end{proposition}
\begin{proof}
整数群的酉表示由一个酉算子 \(U\) 决定。其谱包含于 \(\mathbb T\)：\(|z|>1\) 时用 Neumann 级数，\(|z|<1\) 时写 \(U-zI=U(I-zU^*)\)，仍可求逆。故对 Laurent 多项式 \(p\)，\cref{b2:thm:continuous-functional-calculus}给出
\(\bigl\lVert p(U)\bigr\rVert\le\max_{|z|=1}\bigl|p(z)\bigr|\)。反过来，每个标量 \(z\in\mathbb T\) 都给出一维酉表示，所以最大范数正好为这个上确界。

圆周弧长测度 \(\mu\) 取总质量一。函数 \(z^n\)、\(n\in\mathbb Z\)，两两正交且各有范数一，因为 \(\int_{\mathbb T}z^j\,d\mu(z)=0\) 对非零整数 \(j\) 成立。Laurent 多项式由 Stone--Weierstrass 定理在连续函数中一致稠密，连续函数又在 \(L^2\) 中稠密，所以这些函数构成 \(L^2(\mathbb T)\) 的正交基。因此 \(F\delta_n=z^n\) 从有限支集向量的等距映射延伸为酉同构 \(F:\ell^2(\mathbb Z)\to L^2(\mathbb T)\)。这证明所需的\term[Fourier-Plancherel-dingli]{Fourier--Plancherel 定理}[Fourier--Plancherel theorem]的圆周版本，亦见\cite[Theorem II.10.1.15(iii)]{Blackadar2017OperatorAlgebrasCorrected}。

由 \(F\lambda_1\delta_n=z^{n+1}\)，双边移位在此同构下恰为乘坐标函数。乘连续函数 \(p\) 的算子范数为 \(\max|p|\)：上界直接成立，下界可取最大值附近正测度弧段上的归一化示性函数，使像的范数任意接近最大值。因此约化范数与最大范数一致。最后，Laurent 多项式的一致稠密性使完备化得到 \(C(\mathbb T)\)。
\end{proof}

同一个 Fourier 模型还给出 \(L(\mathbb Z)\cong L^\infty(\mathbb T)\)。证明可由交换子完成。若有界算子 \(T\) 与乘坐标函数及其逆交换，它就与全部 Laurent 多项式、继而与全部连续函数的乘法算子交换。弧长测度的正则性使每个可测集合 \(E\) 的示性函数都可由连续函数 \(f_n\) 在 \(L^2\) 中逼近，且可取 \(0\le f_n\le1\)。记乘函数 \(h\) 的算子为 \(M_h\)。对 \(\xi\in L^\infty(\mathbb T)\subset L^2(\mathbb T)\)，有
\[
\|(M_{f_n}-M_{1_E})\xi\|_2
\le\|\xi\|_\infty\|f_n-1_E\|_2\longrightarrow0.
\]
对一般 \(\xi\in L^2(\mathbb T)\)，先用截断将它逼近为有界函数 \(\xi_k\)。由于 \(\|M_{f_n}-M_{1_E}\|\le1\)，
\[
\|(M_{f_n}-M_{1_E})\xi\|_2
\le\|\xi-\xi_k\|_2
+\|(M_{f_n}-M_{1_E})\xi_k\|_2.
\]
先取 \(k\) 再取 \(n\)，即得 \(M_{f_n}\to M_{1_E}\) 的强算子收敛。于是 \(T\) 也与每个示性函数的乘法算子交换。令 \(g=T1\)，对可测集合 \(E\) 有 \(T1_E=1_Eg\)。于是
\[
\|1_Eg\|_2\le\|T\|\cdot\|1_E\|_2.
\]
若 \(|g|>\|T\|\) 的集合有正测度，取其中 \(|g|\ge\|T\|+\varepsilon\) 的正测度子集，便违反此式。故 \(g\in L^\infty\)，\(T\) 在简单函数及其 \(L^2\) 极限上都是乘 \(g\)。同样的论证表明，全部本质有界函数的乘法代数也是自身的交换子。因此乘连续函数代数的交换子和双交换子均为 \(L^\infty\) 乘法代数；\cref{b2:thm:operator-bicommutant}和 Fourier 同构便给出结论。它的元素按几乎处处相等认同，包含弧段的示性函数；这些通常不是连续函数。于是代数有限和、范数完备化与弱算子闭包确实是三个不同层次。一般交换群的 Fourier 变换与群 \(C^*\)-代数分别见\cite[Theorem II.10.1.15; Proposition II.10.2.7]{Blackadar2017OperatorAlgebrasCorrected}；\(L(G)\) 与 \(L^\infty(\widehat G)\) 的联系及其 Hopf--von Neumann 背景另见同书 II.10.8.13。

\ProseParagraphBreak
一个非交换的例子来自两个生成元上的自由群 \(F_2\)：\(L(F_2)\) 是 \(\mathrm{II}_1\) 因子，具有忠实迹态 \(\tau(T)=\langle\delta_e,T\delta_e\rangle\)，却没有非零极小投影。其因子性与 \(F_2\) 中每个非单位元的共轭类均无限有关：中心元的平方可和系数在共轭类上恒定，因而只能支撑于单位元。迹与投影类型的完整论证见\cite[\nopp III.3.3.2--III.3.3.7]{Blackadar2017OperatorAlgebrasCorrected}。

\subsection{代数张量积与算子张量范数}
\label{b2:subsec:D05:02}

第六章的张量积只有有限和。本小节固定两个非零复 \(C^*\)-代数 \(A,B\)，用 \(A\odot B\) 特别表示它们在复数域上的代数张量积，规定 \((a\otimes b)^*=a^*\otimes b^*\)。符号改变只是为了与下面的完备化区分。
\symidx{\(A\odot B\)：两个复代数的代数张量积}

\ProseParagraphBreak
取 \(A,B\) 的忠实 \(*\)-表示 \(\pi_A:A\to B(H_A)\)、\(\pi_B:B\to B(H_B)\)，先构造表示所作用的空间。在复向量空间的代数张量积 \(H_A\odot H_B\) 上规定
\[
\langle\xi\otimes\eta,\xi'\otimes\eta'\rangle
=\langle\xi,\xi'\rangle\langle\eta,\eta'\rangle,
\]
并对第一变量共轭线性、对第二变量线性地延拓。任意有限张量都落在两侧某个有限维子空间的张量积中；取这两个子空间的正交归一基后，张量积基也正交归一，所得平方范数就是系数绝对值的平方和，所以这个内积正定。按内积范数完备化，得到 Hilbert 张量积 \(H_A\otimes H_B\)。有界算子 \(S\in B(H_A)\)、\(T\in B(H_B)\) 在有限张量上按 \(\xi\otimes\eta\mapsto S\xi\otimes T\eta\) 作用；将有限张量按一侧正交归一基展开，分别作用于两侧，得到范数上界 \(\|S\|\cdot\|T\|\)，故作用连续延伸到这个 Hilbert 张量积。

\ProseParagraphBreak
现在换到代数 \(A\odot B\) 上：令 \(a\otimes b\) 在已经完成的 Hilbert 空间 \(H_A\otimes H_B\) 上按 \(\pi_A(a)\otimes\pi_B(b)\) 作用，并对有限和线性延拓。所得算子范数称为\term[kongjian-zhangliang-fanshu]{空间张量范数}[spatial tensor norm]，亦记 \(\|\cdot\|_{\min}\)。随后按这个范数完备化的是 \(A\odot B\)，得到 \(A\otimes_{\min}B\)。前一次完备化得到表示空间，后一次得到 \(C^*\)-代数；它们作用于不同的对象。

\ProseParagraphBreak
另一方面，另取复 Hilbert 空间 \(H\)，在其上取像互相交换的 \(*\)-表示 \(\pi:A\to B(H)\)、\(\rho:B\to B(H)\)，考察
\[
\sum_i a_i\otimes b_i\longmapsto\sum_i\pi(a_i)\rho(b_i).
\]
记这个诱导的 \(*\)-表示为 \(\pi\cdot\rho\)。对所有这种表示取范数上确界，得到\term[zuidazhangliang-fanshu]{最大张量范数}[maximal tensor norm] \(\|\cdot\|_{\max}\)。\cref{b2:prop:cstar-star-hom-general}保证两个表示均收缩，因而
\[
\left\|\sum_i\pi(a_i)\rho(b_i)\right\|
\le\sum_i\|a_i\|\,\|b_i\|,
\]
所取上确界有限。按这个范数完备化 \(A\odot B\)，所得代数记为 \(A\otimes_{\max}B\)。最小、最大范数比较的对象都是代数 \(A\odot B\)，属于后一次完备化的选择。
\symidx{\(A\otimes_{\min}B\)：空间张量积的范数完备化}
\symidx{\(A\otimes_{\max}B\)：最大张量积的范数完备化}

\ProseParagraphBreak
忠实表示的存在性由\cref{b2:thm:gelfand-naimark-unital,b2:prop:cstar-unitization-representation}保证。矩阵系数能分离有限张量，故 \(\pi_A\otimes\pi_B\) 在 \(A\odot B\) 上单射，空间范数确为范数。两个算子因子的作用彼此交换，所以空间表示也是最大范数定义中容许的一种表示；这同时保证最大范数的非退化性。次乘性与对合的等距性由各表示继承，而
\[
\|x^*x\|_{\max}
=\sup_{\pi,\rho}\|(\pi\cdot\rho)(x)\|^2
=\|x\|_{\max}^2
\]
给出 \(C^*\) 恒等式。

\begin{claim}[最小与最大 \(C^*\) 张量范数]\label{b2:claim:cstar-tensor-extrema}
设 \(A,B\) 为非零复 \(C^*\)-代数，\(A\odot B\) 为其复代数张量积。取忠实表示 \(\pi_A:A\to B(H_A)\)、\(\pi_B:B\to B(H_B)\)，以 \(\pi_A\otimes\pi_B\) 在 \(H_A\otimes H_B\) 上的算子范数定义 \(\|\cdot\|_{\min}\)；对同一 Hilbert 空间上像互相交换的全部 \(*\)-表示 \(\pi,\rho\)，以 \(x=\sum_i a_i\otimes b_i\) 对应算子 \(\sum_i\pi(a_i)\rho(b_i)\) 的范数上确界定义 \(\|x\|_{\max}\)。空间范数不依赖所选的两个忠实表示。对 \(A\odot B\) 上任意 \(C^*\) 范数 \(\gamma\)，都有
\[
\|x\|_{\min}\le\gamma(x)\le\|x\|_{\max}
\qquad(x\in A\odot B),
\]
且 \(\gamma(a\otimes b)=\|a\|\,\|b\|\) 对所有 \(a\in A,b\in B\) 成立。
\end{claim}

这项一般定理见 Blackadar\cite[\nopp II.9.1.3; II.9.5.1--II.9.5.2]{Blackadar2017OperatorAlgebrasCorrected}。最大性的含幺情形可以直接说明：设 \(C=A\otimes_\gamma B\)，在 \(C\) 的忠实 Hilbert 空间表示中，映射 \(a\mapsto a\otimes1_B\)、\(b\mapsto1_A\otimes b\) 是像互相交换的 \(*\)-同态，由\cref{b2:prop:cstar-star-hom-general}自动收缩。因此 \(C\) 的算子范数受最大范数控制，即 \(\gamma\le\|\cdot\|_{\max}\)。

极小性所需的分析工具更多。一个态若不能写成两个不同态的非平凡凸组合，就称为纯态。标准证明先建立纯态与不可约 GNS 表示的对应，再证明纯态乘积能延拓到 \(A\otimes_\gamma B\) 上的态，由这些表示控制空间范数。这里用到态空间的弱星紧性、紧凸集的极点与分离定理，以及交换因子张量积的理想结构；非含幺情形还要用近似单位。相关结果分别见同书 II.6.3--II.6.4、Theorem II.9.4.4 和 II.9.5.1；非含幺表示的处理见 Theorem II.9.2.1 与 Corollary II.9.2.2。这些输入超出此前的 GNS 构造，代数张量积的双线性泛性质本身并不能推出上述范数比较。

在矩阵例子中，\(E_{ij}\otimes E_{ab}\) 送到 \(\mathbb C^n\otimes\mathbb C^m\) 上相应的矩阵单位，给出
\[
M_n(\mathbb C)\odot M_m(\mathbb C)\cong M_{nm}(\mathbb C).
\]
这一代数同构也保持共轭转置；空间范数就是右边的算子范数。最大范数也相同：任何 \(C^*\) 范数 \(\gamma\) 都使这个有限维空间完备，由\cref{b2:thm:continuous-functional-calculus}作用于正矩阵 \(x^*x\)，其范数为最大特征值，所以 \(\gamma(x)^2=\gamma(x^*x)\) 就是通常算子范数的平方。有限维空间已经完备，因此该例不增加新元素。相比之下，\(C\bigl([0,1]\bigr)\odot C\bigl([0,1]\bigr)\) 的元素仅是有限和 \(\sum_i f_i(x)g_i(y)\)；这些函数含常数、分离正方形上的点且对共轭封闭，故由 Stone--Weierstrass 定理在一致范数下稠密于 \(C\bigl([0,1]^2\bigr)\)，但不包含每一个二元连续函数。\cref{b2:ex:D-tensor-completion}给出这样的具体函数及一致逼近，说明完备化在这里有实际作用。

\ProseParagraphBreak

Gelfand 对偶允许把交换函数代数当作拓扑空间的一种代数描述。将其推广为“非交换空间”的研究方法时，需要先指定保留哪些数据：\(C^*\)-代数重视范数与连续表示，von Neumann 代数重视弱闭包与测度结构，附录C的 Azumaya 理论则固定交换基环并研究局部的代数矩阵结构。这些问题有联系，却不能仅凭名称就认作同一种几何。

\ProseParagraphBreak
算子 K 理论通过稳定矩阵中的投影、可逆元及连续同伦研究代数；第四卷附录B的代数 K 理论初步只能提供部分比较语言。Tomita--Takesaki 理论\footnote{\OriginalHanName{冨田稔}（假名：とみた みのる，罗马音：Minoru Tomita，1924-2015），日本数学家，建立 von Neumann 代数的模理论。竹崎正道（假名：たけさき まさみち，罗马音：Masamichi Takesaki，1933-），日本数学家，系统阐述并发展这一理论。}从 von Neumann 代数的态和权构造模自同构群；这里的“模”是该理论的专名，不是本卷任意环模的简称。
